% !TeX root = ../algebralineare.tex

\section{Forme quadratiche} \label{sec:quadform}
\subsection{Definizione di forma quadratica} \label{subsec:quadformdef}
    Un'applicazione $q: \mathbb{R}^{n} \rightarrow \mathbb{R}$ si dice forma quadratica se è determinata da un polinomio di grado 2:
    \[
        \begin{aligned}
        q(x_1, x_2, ..., x_n)
            &= \sum_{i = 1}^{n} \sum_{j = 1}^{n} a_{ij} x_i x_j \\
            &\implies q(x_1, x_2, ..., x_n) = \mathbf{x}^{T} A \mathbf{x}, \\
        A &= \begin{bmatrix}
            a_{11} & a_{12} & ... & a_{1n} \\
            a_{21} & a_{22} & ... & a_{2n} \\
            \vdots & \ddots & \ddots & \vdots \\
            a_{n1} & a_{n2} & ... & a_{nn} \\
        \end{bmatrix},
        &\qquad
        \mathbf{x} &= \begin{bmatrix}
            x_1 \\
            x_2 \\
            \vdots \\
            x_n
        \end{bmatrix}
        \end{aligned}
    \]
    Prendendo per esempio il caso $n = 3$, si ha che:
    \[
        A = \begin{bmatrix}
            a & b & c \\
            d & e & f \\
            g & h & i
        \end{bmatrix}
        \implies
        q(x, y ,z) = \mathbf{x}^{T} A \mathbf{x} = \begin{bmatrix} ax + dy + gz & bx + ey + hz & cx + fy + iz \end{bmatrix} \begin{bmatrix} x \\ y \\ z \end{bmatrix}
    \]
    Arrivando quindi a
    \[
        ax^2 + ey^2 + iz^2 + (b + d)xy + (c + g)xz + (f + h)yz
    \]
    E' così possibile definire la matrice unica (e simmetrica) associata alla forma quadratica:
    \[
        B = \begin{bmatrix}
            a & \frac{(b+d)}{2} & \frac{(c+g)}{2} \\
            \frac{(b+d)}{2} & e & \frac{(f+h)}{2} \\
            \frac{(c+g)}{2} & \frac{(f+h)}{2} & i
        \end{bmatrix}
        \implies
        B = \frac{1}{2} (A + A^T)
    \]



\vspace{10pt}
\subsection{Cambiamento di variabile} \label{subsec:quadformchange}
    Sia $q: \mathbb{R}^{n} \rightarrow \mathbb{R}$ una forma quadratica con forma ridotta $q(\mathbf{x}) = \mathbf{x}^{T} B \mathbf{x}$. E' possibile allora passare dalla variable $\mathbf{y} \in \mathbb{R}^{n}$ alla variabile $\mathbf{x} \in \mathbb{R}^{n}$
    attraverso una matrice di cambio di base $P^{F, E}$ (dove $E$ ed $F$ sono, rispettivamente, le basi di $\mathbf{x}$ e $\mathbf{y}$). Dunque:
    \[
        q(\mathbf{x}) = q(P \mathbf{y}) = (P \mathbf{y})^{T} B (P \mathbf{y}) = \mathbf{y}^T (P^T B P) \mathbf{y} = r(\mathbf{y})
    \]
    Ora, in virtù del fatto che $B$ è simmetrica, per il teorema spettrale, si ha che esiste un matrice ortogonale di autovettori $P$ tale per cui $D = P^{T} B P$, dove $D$ è la matrice diagonale di autovalori. Quindi, operando un cambio di variable secondo $P$
    \[
        q(\mathbf{x}) = r(\mathbf{y}) = \mathbf{y}^{T} (P^T B P) \mathbf{y} = \mathbf{y}^{T} D \mathbf{y} = \sum_{i = 1}^{n} \lambda_i y_i^2  \quad \textnormal{con} \: \mathbf{y} = P^{T} \mathbf{x}
    \]
    La trasformazione di una forma quadratica in forma diagonale è chiamata trasformazione di Jacobi.



\vspace{10pt}
\subsection{Segno di una forma quadratica} \label{subsec:quadformsign}
    E' possibile classificare le forme quadratiche in base al segno della loro immagine.
    \begin{enumerate}[label=(\arabic*)]
        \item Definita positiva: $q(\mathbf{x}) > 0$ per ogni $\mathbf{x} \in \mathbb{R}^{n}$
        \item Semidefinita positiva: $q(\mathbf{x}) \geq 0$ per ogni $\mathbf{x} \in \mathbb{R}^{n}$
        \item Definita negativa: $q(\mathbf{x}) < 0$ per ogni $\mathbf{x} \in \mathbb{R}^{n}$
        \item Semidefinita negativa: $q(\mathbf{x}) \leq 0$ per ogni $\mathbf{x} \in \mathbb{R}^{n}$
        \item Non definita: $\exists \mathbf{x}_{1}, \mathbf{x}_{2} \in \mathbb{R}^{n} \; | \; q(\mathbf{x}_{1}) > 0 \; \land \; q(\mathbf{x}_{2}) < 0$
    \end{enumerate}

    Grazie alla trasformazione di Jacobi [\Cref{subsec:quadformchange}], è possibile tramutare una generica forma quadratica $q(\mathbf{x})$ in forma ``diagonale'' $r(\mathbf{y}) = \lambda_1 y_{1}^{2} + ... + \lambda_n y_{n}^{2}$ dove $\mathbf{y} = P^{T} \mathbf{x}$ e $P$ è una matrice ortogonale di autovettori di $B$.
    Di conseguenza, è possibile affermare che:
    \begin{enumerate}[label=(\arabic*)]
        \item $q(\mathbf{x})$ è definita positiva $\iff$ $\lambda_1, ..., \lambda_n > 0$
        \item $q(\mathbf{x})$ è semidefinita positiva $\iff$ $\lambda_1, ..., \lambda_n \geq 0$
        \item $q(\mathbf{x})$ è definita negativa $\iff$ $\lambda_1, ..., \lambda_n < 0$
        \item $q(\mathbf{x})$ è semidefinita negativa $\iff$ $\lambda_1, ..., \lambda_n \leq 0$
        \item $q(\mathbf{x})$ è non definita $\iff$ $\exists \lambda_1, \lambda_2 \in \Lambda_B \; | \; \lambda_1 < 0 \; \land \; \lambda_2 > 0$
    \end{enumerate}



\vspace{10pt}
\subsection{Coniche} \label{subsec:quadformconic}
    Le coniche sono varietà algebriche affini monodimensionali o zerodimensionali definite in $\mathbb{R}^2$ da un polinomio di secondo grado della forma:
    \[
        q(x, y) = A x^2 + B xy + C y^2 + D x + E y + F = 0
        \Rightarrow
        q(x, y) = \begin{bmatrix} x & y\end{bmatrix} \begin{bmatrix} A & B/2 \\ B/2 & C \end{bmatrix} \begin{bmatrix} x \\ y \end{bmatrix} + \begin{bmatrix} x & y \end{bmatrix} \begin{bmatrix} D \\ E \end{bmatrix}+ F = 0
    \]
    Scritta in forma quadratica omogenea si ha:
    \[
        q(x, y) = \begin{bmatrix} x & y & 1 \end{bmatrix} \begin{bmatrix} A & B/2 & D/2 \\ B/2 & C & E/2 \\ D/2 & E/2 & F \end{bmatrix} \begin{bmatrix} x \\ y \\ 1 \end{bmatrix}, \quad A_q = \begin{bmatrix} A & B/2 \\ B/2 & C \end{bmatrix}, \; B_q = \begin{bmatrix} A & B/2 & D/2 \\ B/2 & C & E/2 \\ D/2 & E/2 & F \end{bmatrix}
    \]
    La matrice $A_q$ è chiamata matrice della forma quadratica, mentre la matrice $B_q$ è chiamata matrice dell'equazione quadratica.

    \vspace{10pt}
    \subsubsection{Forma canonica di una conica} \label{subsubsec:quadformcanonicconic}
        La conica si dice invece in forma ``canonica'' quando $B_q$ è diagonale (dato che $A_q$ è uguale a $B_q$ senza la terza riga e terza colonna, se $B_q$ è diagonale, lo è anche $A_q$) oppure, in altre parole, quando l'equazione della conica diventa:
        \[
            A' x^2 + B' y^2 + C' = 0
        \]

        \vspace{10pt}
        Data una generica conica in forma
        \[
            \begin{bmatrix} x & y\end{bmatrix} \begin{bmatrix} A & B/2 \\ B/2 & C \end{bmatrix} \begin{bmatrix} x \\ y \end{bmatrix} + \begin{bmatrix} x & y \end{bmatrix} \begin{bmatrix} D \\ E \end{bmatrix} + F = 0
        \]
        E' sempre possibile trasformarla in forma canonica secondo un sistema di coordinate ruotato e traslato. Dato che la matrice $A_q$ è simmetrica è possibile, per il teorema spettrale, diagonalizzarla attraverso una matrice ortogonale $P$ di autovettori di $A_q$. Si considerino dunque $x,\: y$ come il risultato di un cambio di variabile $\begin{bmatrix} x \\ y \end{bmatrix} = P \begin{bmatrix} x' \\ y' \end{bmatrix}$.
        Allora l'equazione diventa:
        \[
            \begin{bmatrix} x' & y' \end{bmatrix} D \begin{bmatrix} x' \\ y' \end{bmatrix} + \begin{bmatrix} x' & y' \end{bmatrix} P^{T} \begin{bmatrix} D \\ E \end{bmatrix} + F = 0 \implies \lambda_1 x'^2 + \lambda_2 y'^2 + D' x' + E' y' + F = 0
        \]
        Ciò significa che nel sistema di rifermento ruotato $x', \: y'$ la conica è data dai punti che soddisfano l'equazione soprastante. Per trovare la forma canonica, bisogna trovare un nuovo sistema di riferimento $x'', \: y''$ dove la conica è data dai punti che soddisfano un'equazione canonica.
        Assumendo che $\lambda_1$ e $\lambda_2$ siano diversi da 0, si consideri:
        \[
            \begin{aligned}
                \lambda_1 \left( x'^2 + \frac{D'x'}{\lambda_1} \right) &= \lambda_1 \left[\left(x' + \frac{D'}{2 \lambda_1}\right)^2 - \left( \frac{D'}{2 \lambda_1}\right)^2 \right]  \\
                \lambda_2 \left( y'^2 + \frac{E'y'}{\lambda_2} \right) &= \lambda_2 \left[\left(y' + \frac{E'}{2 \lambda_2}\right)^2 - \left( \frac{E'}{2 \lambda_2}\right)^2 \right]
            \end{aligned}
        \]
        Da cui si ottiene
        \[
            \lambda_1 \left(x' + \frac{D'}{2 \lambda_1}\right)^2 - \frac{D'^2}{4 \lambda_1} + \lambda_2 \left(y' + \frac{E'}{2 \lambda_2}\right)^2 - \frac{E'^2}{4 \lambda_2} + F = 0
        \]
        Sia allora $x' + D '/ 2 \lambda_1 = x''$ e $y' + E' / 2 \lambda_2 = y''$, con $g = F - \frac{D'^2}{4 \lambda_1} - \frac{E'^2}{4 \lambda_2}$. Si arriva infine a
        \[
            \lambda_1 x'' + \lambda_2 y'' + g = 0
        \]
        Quindi, ripercorrendo dall'inizio, si osserva che il nuovo sistema di riferimento $x', \: y'$ è dato da una traslazione del sistema $x'', \: y''$ $\left( \begin{bmatrix} x' \\ y' \end{bmatrix} = \begin{bmatrix} x'' \\ y'' \end{bmatrix} + \begin{bmatrix} - D' / 2 \lambda_1 \\ - E' / 2 \lambda_2 \end{bmatrix} \right)$. Le coordinate di questa traslazione sono chiamate coordinate del centro della conica.
        Tenendo conto che $\begin{bmatrix} - D' / 2 \lambda_1 \\ - E' / 2 \lambda_2 \end{bmatrix}$ sono le coordinate del centro della conica nel sistema di riferimento ruotato (quindi $\begin{bmatrix} - D' / 2 \lambda_1 \\ - E' / 2 \lambda_2 \end{bmatrix} = P^{T} \begin{bmatrix} x_0 \\ y_0 \end{bmatrix}$), si ha che:
        \[
            \begin{bmatrix} x \\ y \end{bmatrix} = P \begin{bmatrix} x' \\ y' \end{bmatrix} = P \left( \begin{bmatrix} x'' \\ y'' \end{bmatrix} + \begin{bmatrix} - D' / 2 \lambda_1 \\ - E' / 2 \lambda_2 \end{bmatrix} \right) = P \left( \begin{bmatrix} x'' \\ y'' \end{bmatrix} + P^{T} \begin{bmatrix} x_0 \\ y_0 \end{bmatrix} \right) = P \begin{bmatrix} x'' \\ y'' \end{bmatrix} + \begin{bmatrix} x_0 \\ y_0 \end{bmatrix}
        \]
        In altre parole, la conica originale in $x, \: y$ non è che il risultato di due trasformazioni, più precisamente una rotazione ed una traslazione, applicate ad una conica canonica.

    \vspace{10pt}
    \subsubsection{Tipi di coniche} \label{subsubsec:quadformconictypes}
        Ora, a seconda dei valori di $\det(B_q)$ e $\det(A_q)$ è possibile determinare il tipo di conica:
        \begin{enumerate}[label=(\arabic*)]
            \item  $\det(B_q) = 0 \implies$ conica degenere (coppia di rette)
            \item $\det(B_q) \neq 0$  $\det(A_q) > 0 \implies$  ellisse
            \item $\det(B_q) \neq 0$  $\det(A_q) = 0 \implies$  parabola
            \item $\det(B_q) \neq 0$  $\det(A_q) < 0 \implies$  iperbole
        \end{enumerate}
        Nel caso dell'ellise, se $A = C$ e $B = 0$, allora la conica è un cerchio.
